\documentclass[a4paper,14pt]{extarticle}
\usepackage{styles}

\begin{document}

\litSubsubsection{2.1.1 Genetic Algorithms Parallelization Potential}

Genetic algorithms (GAs) are metaheuristic optimization techniques inspired by natural selection, which simulate evolutionary processes through key operations: fitness evaluation, selection, crossover, and mutation. GAs are stochastic global search methods that combine two strategies: exploiting better solutions and exploring the search space~\cite{cheng2020parallel}. A GA maintains a population of candidate solutions (individuals), each encoded as a chromosome, and iteratively evolves them over generations. GAs belong to the class of black-box optimization methods where only objective function evaluations are required - no gradient information or problem-specific knowledge is needed~\cite{katoch2021review}. The concept originates from John Henry Holland's work at the University of Michigan~\cite{holland1975adaptation}, whose ``original goal was not to design algorithms to solve specific problems, but rather to formally study the phenomenon of adaptation as it occurs in nature''~\cite{mitchell1998introduction}.

\parahead{Applications and Problem Domains.}
GAs are effective for solving NP-hard problems such as the Traveling Salesman Problem (TSP), Vehicle Routing Problem (VRP), job-shop scheduling, and multi-objective optimization~\cite{goldberg1989genetic,katoch2021review}. They avoid local optima traps that complicate gradient-based methods and are suited for vast, discontinuous search spaces.

\parahead{Core Operations of Genetic Algorithms.}
The GA lifecycle comprises four operations:
\gaop{Fitness Evaluation} computes an objective function value for each individual independently - typically the main computational bottleneck.
\gaop{Selection} chooses parents probabilistically based on fitness (e.g., tournament, roulette wheel).
\gaop{Crossover} combines parent chromosomes to produce offspring (e.g., single-point, uniform).
\gaop{Mutation} introduces random changes to maintain diversity (e.g., bit-flip, swap).

Each generation, $n$ offspring are produced through selection, crossover ($P_c$), and mutation ($P_m$), then replace the current population~\cite{katoch2021review}. Figure~\ref{fig:ga-lifecycle} illustrates this generational loop.

\begin{figure}[H]
\centering
\includegraphics[width=0.45\textwidth]{diagrams/sec2/ga-lifecycle.png}
\caption{GA generational loop. Fitness evaluation is the primary GPU acceleration target due to its embarrassingly parallel nature.}
\label{fig:ga-lifecycle}
\end{figure}

\subsubsection*{Black-Box Optimization Nature and Parallelization Potential}

GAs are black-box optimizers: they evaluate solutions solely through fitness function calls, requiring no gradient computation or analytical models~\cite{katoch2021review}. ``The GA does not require derivative information or any other auxiliary knowledge about the function being optimized''~\cite{mitchell1998introduction}. In formal terms, the objective function is treated as an opaque oracle returning $f(\vect{x})$ given input $\vect{x}$, without revealing gradients or analytical form~\cite{audet2017derivative,nocedal2006numerical}.

This black-box property directly enables parallelization: since fitness evaluation has no inter-individual dependencies, each individual can be evaluated in complete isolation. This transforms population-wide fitness evaluation into an embarrassingly parallel workload, the ideal case for GPU acceleration where thousands of threads execute identical computations on different individuals~\cite{cheng2020parallel}. As Cheng and Gen note, ``on GPU architecture, you can run it with more than hundreds of thousands of massive threads simultaneously''~\cite{cheng2020parallel}.

Fitness evaluation dominates sequential GA execution time: ``the bottleneck of the GAs is the fitness evaluation task''~\cite{salami2002fitness}, and can consume the majority of total computation time for complex fitness functions such as neural network training~\cite{salami2002fitness}.

\end{document}
